\section{Связанные работы}
\label{sec:related}

\textbf{Численная устойчивость и неассоциативность.}
Классический аппарат вычислительной погрешности с плавающей точкой~--- обратный
анализ ошибок Уилкинсона~\cite{Wilkinson1963}, систематизация свойств
IEEE~754 Голдбергом~\cite{Goldberg1991} и оценки погрешности матричных
алгоритмов Хайэма~\cite{Higham2002}~--- задаёт теоретическую рамку работы.
Неассоциативность \texttt{FP32}-суммирования определяет принципиальную
возможность расхождений между реализациями GEMM с различным порядком
накопления; детерминированная верхняя оценка Хайэма (линейный рост ошибки по
размерности редукции $K$) служит ориентиром, достижимым лишь в наихудшем случае.

\textbf{Аппаратные расхождения GPU-реализаций GEMM.}
Побитово-точные модели тензорных ядер~\cite{MMASim2025,Khattak2025} показывают,
что матричные умножители ускорителей не соответствуют IEEE~754 на уровне
инструкции и вносят систематический сдвиг результата; ковариационный
анализ~\cite{Yashwanth2025} опровергает гипотезу о независимости ошибок
округления, выявляя их структурированность. Эти результаты предсказывают
конструктивное накопление сдвига при прохождении через слои.

\textbf{Допустимые уровни погрешности и тензорные компиляторы.}
Работы по обучению в смешанной разрядности~\cite{Micikevicius2018,Wang2018FP8}
устанавливают эмпирические пороги относительного отклонения, при которых
точность перестаёт влиять на сходимость. Самостоятельным источником расхождений
служит компилятор: слияние операций в XLA~\cite{Snider2023}, автоматический
выбор расписания в TVM~\cite{Chen2018TVM} и многоуровневые преобразования
MLIR~\cite{Lattner2021MLIR} изменяют порядок округлений, причём сильнее, чем
смена аппаратной реализации,~--- откуда вытекает требование изолировать
компиляторный и аппаратный вклады.

\textbf{Симуляторы ускорителей и воспроизводимость.}
Систолический массив TPU~\cite{Jouppi2017} реализует уникальную схему
накопления, отличную от CPU и GPU; поведенческие
симуляторы~\cite{Samajdar2020,Rashidi2020} оценивают производительность, но
абстрагируются от побитовой арифметики, тогда как исследование аппаратных
расхождений требует функционального симулятора с точной арифметикой.
Расхождения, изучаемые здесь, воспроизводимы и стационарны при фиксированной
конфигурации компилятора~--- в отличие от недетерминизма атомарных
редукций~\cite{Langroudi2024,ZhangDeterministic2025}.

\textbf{Архитектурная устойчивость сетей.}
Расположение нормализации (Pre-LN или Post-LN) определяет устойчивость градиентов
при обучении~\cite{Xiong2020}; DeepNet~\cite{Wang2022} делает Post-LN пригодным
для глубоких сетей. Те же свойства нормализационного оператора, что стабилизируют
градиенты, влияют и на накопление аппаратных расхождений: нормализационный слой
действует как сжимающее отображение по спектральной норме и выступает стоком
расхождений (раздел~\ref{sec:model}).

\textbf{Место настоящей работы.}
Перечисленные направления рассматривают источники расхождений по отдельности:
арифметику GEMM~\cite{Higham2002,MMASim2025,Khattak2025}, компилятор~\cite{Snider2023},
архитектуру сети~\cite{Xiong2020}. Однако в реальном конвейере симуляции эти
вклады действуют совместно, и их суммарный эффект на послойное накопление
не измерялся: существующие работы либо дают верхние оценки для одной операции,
либо характеризуют аппаратный сдвиг вне контекста сети, либо анализируют
устойчивость архитектуры к возмущениям без привязки к конкретному источнику
вычислительной ошибки. Настоящая работа закрывает этот пробел: предложенный
измерительный конвейер \emph{раздельно} оценивает вклад компилятора и реализации
GEMM в пределах единой цепочки исполнения и прослеживает их распространение по
слоям сети при прямом и обратном проходах, что позволяет установить иерархию
источников и обосновать выбор реализации GEMM в функциональном симуляторе.
